% ============================================================================
% Section 7.1 — Parts of a Circle
% CCSS 7.G.B.4
% ============================================================================
\section{Parts of a Circle}

\begin{stepGoal}
\begin{itemize}
    \item Identify the center, radius, diameter, and circumference of a circle
    \item Use the relationship $d = 2r$ to convert between radius and diameter
\end{itemize}
\end{stepGoal}

\begin{stepVocabBox}
    \vocabItem{Radius}{Distance from the center to any point on the circle.}
    \vocabItem{Diameter}{Distance across the circle through the center; $d = 2r$.}
    \vocabItem{Circumference}{The distance around the outside of the circle.}
\end{stepVocabBox}

\begin{stepsCard}[How to Identify and Use Circle Parts]
    \stepItem{Find the \textbf{center} — the middle point of the circle.}
    \stepItem{Identify the \textbf{radius} (center to edge) or \textbf{diameter} (edge to edge through center).}
    \stepItem{Convert: $d = 2r$ or $r = d \div 2$.}
\end{stepsCard}

% --- Worked Example 1 ---
\begin{stepExample}{A circle has a radius of $6$ cm. Find the diameter.}
    \stepShow{1} The center is the middle point of the circle.

    \stepShow{2} The radius is $6$ cm (from center to edge).

    \stepShow{3} Diameter $= 2r = 2 \times 6 = 12$ cm.
    \stepResult{The diameter is $12$ cm.}
\end{stepExample}

% --- Worked Example 2 ---
\begin{stepExample}{A dinner plate has a diameter of $10$ inches. What is the radius?}
    \stepShow{1} The center is the middle of the plate.

    \stepShow{2} The diameter is $10$ in (across through the center).

    \stepShow{3} Radius $= d \div 2 = 10 \div 2 = 5$ in.
    \stepResult{The radius is $5$ inches.}
\end{stepExample}

\mascotStepTip{The diameter always passes through the center. It is the longest distance across a circle!}

% ============================================================================
% PRACTICE
% ============================================================================
\newpage
\begin{stepPractice}{Parts of a Circle Practice}
    \resetProblems

    \stepPracticeHeader[step1Color]{Identify the Part}
    \begin{multicols}{2}
    \prob The segment from the center to the edge is called the \answerBlank[2.5cm]
    \answer{radius}
    \prob The segment across the circle through the center is called the \answerBlank[2.5cm]
    \answer{diameter}
    \end{multicols}

    \stepPracticeHeader[step2Color]{Find the Diameter}
    \begin{multicols}{2}
    \prob Radius $= 9$ m. Diameter $=$ \answerBlank[2cm]
    \answer{$18$ m}
    \prob Radius $= 3.5$ in. Diameter $=$ \answerBlank[2cm]
    \answer{$7$ in}
    \end{multicols}

    \stepPracticeHeader[step3Color]{Find the Radius}
    \prob Diameter $= 14$ ft. Radius $=$ \answerBlank[2cm]
    \answer{$7$ ft}

    \wordProblem{A bicycle wheel has a diameter of $26$ inches. What is the radius of the wheel?}{inches}
    \answer{$13$ inches}
\end{stepPractice}
